\documentclass[11pt,a4paper]{article}
\usepackage[T1]{fontenc}
\usepackage[utf8]{inputenc}
\usepackage{lmodern,microtype}
\usepackage[margin=28mm]{geometry}
\usepackage{amsmath,amssymb}
\usepackage{booktabs,array,longtable}
\usepackage{enumitem}
\usepackage[dvipsnames]{xcolor}
\usepackage[colorlinks,linkcolor=NavyBlue,citecolor=NavyBlue,urlcolor=NavyBlue]{hyperref}
\linespread{1.06}
\setlength{\emergencystretch}{3em}
%% generated by scripts/process-record.py -- do not edit
\newcommand{\NCommits}{116}
\newcommand{\FirstDay}{2026-09-08}
\newcommand{\LastDay}{2026-09-14}
\newcommand{\NDays}{6}
\newcommand{\NFallacies}{41}
\newcommand{\NErrata}{34}
\newcommand{\NScriptsTS}{38}
\newcommand{\NScriptsPy}{29}
\newcommand{\NData}{470}
\newcommand{\NReviews}{1}
\newcommand{\NTexLines}{5654}
\newcommand{\RepoRev}{238bb3b}
\newcommand{\PagesI}{31}
\newcommand{\PagesII}{25}
\newcommand{\PagesIII}{27}
\newcommand{\PagesIV}{33}
\newcommand{\CommitsTable}{\begin{tabular}{lr}\toprule day & commits\\\midrule
2026-09-08 & 1\\
2026-09-09 & 8\\
2026-09-10 & 39\\
2026-09-11 & 14\\
2026-09-13 & 14\\
2026-09-14 & 40\\
\bottomrule\end{tabular}}


% ---- scaffold macros: visible in the draft, silenced by \finalversion
\newif\iffinal
%\finaltrue
\newcommand{\interview}[2]{\iffinal\else\par\noindent\colorbox{Goldenrod!25}{\parbox{\dimexpr\linewidth-2\fboxsep}{\small\textbf{Interview Q#1.} #2}}\par\fi}
\newcommand{\todo}[1]{\iffinal\else\par\noindent\colorbox{Red!12}{\parbox{\dimexpr\linewidth-2\fboxsep}{\small\textbf{To do.} #1}}\par\fi}
\newcommand{\evidence}[1]{\iffinal\else{\footnotesize\textcolor{Gray}{[evidence: #1]}}\fi}
\newcommand{\file}[1]{\texttt{#1}}

\title{Four papers in six days:\\ a case study of AI-driven mathematical research\\ under a verifiable process}
\author{Jasper Reichardt}
\date{Draft scaffold, 14 September 2026}

\begin{document}
\maketitle

\begin{abstract}
\noindent
Between \FirstDay{} and \LastDay{} a software engineer with no research background in analytic number theory
produced, with a large language model doing almost all of the mathematical, computational and textual work,
a series of four research papers on the pair singular series of prime values of polynomials
(\PagesI{}, \PagesII{}, \PagesIII{} and \PagesIV{} pages). The mathematics is not the subject of this paper. Its subject is
the \emph{process}: how a test-driven engineering habit and a small set of structural guardrails (a public
knowledge base of state and errors, a project-wide errata file, adversarial readings by fresh model instances,
external assessments archived against the revision they read, a two-routes rule for every constant, and a
stopping criterion) turned an inherently unreliable generator of mathematics into a process whose output can be
checked, and was checked, by people who did not take part in it. We document the set-up, the evolution of the
prompting, the guardrails and when each was introduced, the complete record of \NFallacies{} recorded
fallacies and \NErrata{} errata items, who caught which class of error, and what the process could not catch. The
programme's mathematical status is reported exactly as its own errata report it: theorems proved and read,
theorems retracted, and problems left open.
\todo{Rewrite the abstract once the interview is in; the numbers are generated from the repository by
\file{scripts/process-record.py}.}
\end{abstract}

\tableofcontents

\section{Introduction}\label{sec:intro}

\paragraph{The question.} Large language models can write mathematics that looks right. The question this case
study addresses is narrower and more useful than whether they can do research: under what \emph{process} does a
non-specialist, working with such a model, produce mathematical output whose correctness does not depend on
trusting either of them? The programme documented here is one data point, with a complete public record, and
this paper is the reading of that record.

\paragraph{What is claimed and what is not.} We claim that the process described in \S\ref{sec:guardrails} produced
(i) a body of work in which every statement is labelled proved, conditional, conjectural or open, (ii) a record of
every error found, by whom, and how it was fixed, and (iii) numbers that can be regenerated from the repository. We
do \emph{not} claim that the mathematics is important, nor that it is correct beyond what its errata say; as of
this writing the four papers have been read adversarially by model instances without session context and by two
outside readers, and not yet by a specialist in the field (\S\ref{sec:limits}). The case study has $n=1$.

\paragraph{The subject programme in one paragraph.} For an irreducible polynomial $f$, the pair singular series
$S_f(h)$ is the explicit weight that, under the Hardy--Littlewood conjecture, governs simultaneous primality of
$f(t)$ and $f(t+h)$. Paper~I proves an exact identity for $\sum_{h\le H}(S_f(h)-C(f)^2)$, splits it into a diagonal
whose Dirichlet series is $\zeta_K(s+1)E_f(s)$ and an off-diagonal, and states a conjecture in Ces\`aro form.
Paper~II identifies $E_f$ as an infinite product of Artin $L$-functions and proves a conditional explicit formula.
Paper~III decomposes the off-diagonal into pieces and proves everything except one window. Paper~IV finds, and
partly proves, that the pieces oscillate at the spectral parameters of even Maass forms. Nothing in the four papers
proves anything about primes without Hardy--Littlewood, and nothing bears on the Riemann Hypothesis.
\evidence{\file{research/KNOWLEDGE.md} \S0; root \file{README.md}}

\paragraph{Reader's guide.} \S\ref{sec:setup} describes the author, the tools and the timeline. \S\ref{sec:origin}
describes how a visualisation tool with a test suite became a research programme. \S\ref{sec:guardrails} lists the
guardrails and their origin. \S\ref{sec:prompting} traces the evolution of the prompting. \S\ref{sec:errors} reads
the error record. \S\ref{sec:verification} describes what was verified and how. \S\ref{sec:limits} says what the
process cannot do. \S\ref{sec:recommend} draws recommendations. Appendix~\ref{app:interview} contains the interview,
Appendix~\ref{app:record} the process record.

\section{Set-up}\label{sec:setup}

\subsection{The author}
The author is a freelance software engineer, entrepreneur and trader running several projects at once, with a
bachelor's degree in software engineering. The hardest mathematics course of that degree was stochastics; Fourier
transforms and Markov chains were fluent at the time, and professional work had included specialised algorithms
(for instance enumerating all cycles in a weighted bidirectional graph up to a given size). No analytic number
theory. What the author could read in a number-theory paper on 8~September was the prose, the meaning of an
asymptotic statement and why it matters, and concepts such as the Bateman--Horn constant $C(f)$ and the comparison of
polynomials by their arithmetic; not the proofs. The curiosity about primes was old, and the prime-rich diagonals of
the Ulam spiral were the visual intuition from which the study started. \evidence{Interview Q1}
\interview{1b}{Still open from Q1: prior experience with LLMs in engineering work before this project (how long,
which tools, in what role: autocomplete, pair programming, agents?).}

\subsection{The tools}
The work was done in a terminal agent (Claude Code) driving Claude models: Claude Opus~5 for most of the
programme, Claude Fable~5.1 for the correction phase from 13~September. Fresh model instances without access to
the working session were used as adversarial readers (\S\ref{sec:guardrails}); a cheaper model was used for
simpler tasks after rate limits made this necessary. One external assessment was produced by a different vendor's
model (ChatGPT) on request of the author and archived as a PDF against the revision it read; a second outside
screening was produced by a separate Claude instance without repository access beyond the papers.
\evidence{\file{research/reviews/2026-09-12-assessment-a372a63.pdf}; ERRATA 11--15, 26; disclosure paragraphs in each paper}
Software: TypeScript (\file{tsx}) for sieves and exact rational arithmetic, Python with \file{mpmath} for high
precision, \file{vitest} for the workbench tests, \LaTeX{} via \file{tectonic}, and downloaded Maass-form
coefficients from the LMFDB. Every number in the four papers is produced by a script in the repository.
\interview{2}{Tools and cost: sessions per day, hours per day, approximate token or money cost, other tools
(ChatGPT for the assessment, anything else), and whether anything was done outside the agent.}

\subsection{The timeline}
\begin{center}\small\CommitsTable\end{center}
\noindent\NCommits{} commits over \NDays{} days; \NTexLines{} lines of \LaTeX{} across the four papers;
\NScriptsTS{} TypeScript and \NScriptsPy{} Python sources; \NData{} data files; \NReviews{} archived external
assessment. The phases, from the commit log and the knowledge base:
\begin{description}[leftmargin=2.2em,itemsep=1pt]
\item[8 Sep] The Ulam-nD workbench (an $n$-dimensional prime-spiral explorer) is built with a test suite; the
  \emph{thesis suite} (\file{research/experiments/run.ts}) states the author's four claims as pass/fail tests and
  refutes three of them the same evening. The one surviving lead is a sub-Poisson variance along polynomial rays.
\item[9 Sep] Paper~I (versions 6--7) and the scaffold of paper~II; the knowledge base with timeline, status board,
  fallacies and rated paths is started; the first ``second opinion'' by another model instance is assessed
  line by line.
\item[10 Sep] Paper~III written and self-reviewed (39 commits in one day); paper~IV scaffolded, and the ``second
  spectrum'' (Maass-form oscillations in the pieces) found numerically at 21:02.
\item[11 Sep] Authorship decision: the model is removed as co-author, a uniform disclosure section replaces it
  (ERRATA~8). Paper~IV's theorem written.
\item[12 Sep] The external assessment (three findings, all confirmed); the obstruction in paper~IV identified
  exactly; the first error in paper~IV's theorem found by the assessment.
\item[13 Sep] Corrections applied (ERRATA 11--23); paper~IV repaired at level $\Gamma_0(e)$; the remainder bound
  found not to close and said so; the smooth-window theorem proved and confirmed on three spectral lines; six
  adversarial readings in one day, errors in five of them.
\item[14 Sep] The last two internal readings applied (ERRATA 24--25); the internal phase declared finished by the
  criterion set the day before; a second outside reader finds an unproved non-vanishing claim in paper~I
  (ERRATA~26), fixed the same day.
\end{description}
\evidence{\texttt{git log}; \file{KNOWLEDGE.md} \S1 Timeline}

\section{From a visualisation tool to a research programme}\label{sec:origin}

The programme did not start as one. It started as a web application for looking at prime spirals in higher
dimensions, built to test a private thesis (that Ulam-spiral pseudo-patterns become polynomial patterns in higher
dimensions, that Fibonacci dimensions are special, and that this might say something about the Riemann
Hypothesis). Two engineering decisions on the first day shaped everything after:
\begin{enumerate}[itemsep=1pt]
\item \textbf{The thesis was written as a test suite.} Each claim became a hypothesis with a pass/fail criterion,
  a deterministic seed, a raw report (\file{REPORT.md}) and a separate reading (\file{CONCLUSIONS.md}). Three of the
  four claims failed the tests within hours; the tool had been built to falsify, and it did.
  \evidence{\file{research/experiments/CONCLUSIONS.md}}
\item \textbf{The honest framing was agreed inside the tool.} The theory panel of the application stated that
  spiral patterns are a Bateman--Horn main-term phenomenon, that the Riemann Hypothesis is an error-term
  statement, and that no link between them is known. This sentence survived into every paper.
\end{enumerate}
The one lead the tests left standing, a sub-Poisson variance along polynomial rays, is the Montgomery--Soundararajan
second moment for polynomial prime values; its diagonal weight $\omega_f(p)/(p-\omega_f(p))$ is where paper~I began.
\interview{3}{The turn: when and why did you decide the tool's residue was worth a paper, and then four? What did you
expect at that point, and what did the model tell you to expect?}

\section{The guardrails}\label{sec:guardrails}

Each guardrail below is stated with what it prevents, when it entered the process, and, where the record shows it,
what triggered it. The column ``origin'' is to be completed from the interview.

\begin{longtable}{@{}p{3.6cm}p{5.6cm}p{2.2cm}p{3.6cm}@{}}
\toprule
guardrail & what it does & since & origin / trigger\\
\midrule
\endhead
Knowledge base (\file{KNOWLEDGE.md}): one-paragraph state, ``where things stand and what to do next'', status board,
timeline & Forces every session to start from the recorded state, not from the model's memory of it; makes
``what is proved'' a maintained list rather than a feeling & 9 Sep & \interview{4a}{whose idea; what went wrong
before it existed}\\
Fallacy list (\S3 of the knowledge base, F1--F\NFallacies) & Every error is written down with its fix and a rule
that would have prevented it; read before repeating & 9 Sep & \\
Project-wide errata (\file{ERRATA.md}), never deleted & Corrections are additive and dated; a reader can see what
each version claimed & 9 Sep (per paper), 13 Sep (project-wide) & \\
Every number from a script in the repository & No number in a paper is typed by hand or by the model & 8 Sep &
inherited from the workbench's macro generation\\
Two independent routes for every constant (F9) & Product vs.\ $L$-values; partial sum plus analytic tail vs.\ exact
& 9--10 Sep & a truncated Euler product multiplying $H$ (F4)\\
``Difference it first'' (F26); recompute every pole order in a ``same proof'' generalisation (F27); list every
summation condition and check group invariance (F28) & Rules extracted from the three findings of the external
assessment & 13 Sep & the assessment\\
Adversarial reading by a fresh model instance with no session context & Separates the generator from the checker;
a reader who did not write the proof and cannot see the session & 9 Sep (second opinion), systematic from 12 Sep &
\\
External assessment archived with the revision hash & The review and the text it reviewed stay paired & 12 Sep & \\
Disclosure paragraph, identical in all papers; single human author & Publisher policy and responsibility & 11 Sep
(ERRATA~8) & author's decision\\
``Correct or leave open; never patch with a weaker claim about a different object'' & Decides what a repair may do
& 13 Sep & the author, choosing between repair routes for paper~IV\\
Stopping criterion for the internal phase & Ends open-ended polishing: every statement proved and read at least
once, or labelled open; every number reproducible; errata current; open problems stated with obstructions & 13 Sep &
the author asking ``are we over-rigoring?''\\
``Do not claim'' list in the knowledge base & Sentences the model is forbidden to write (e.g.\ anything about
primes without Hardy--Littlewood, anything about RH) & 9 Sep & \\
\bottomrule
\end{longtable}
\evidence{\file{KNOWLEDGE.md} \S0b, \S3, \S6; \file{ERRATA.md}}
\interview{4}{For each row: was it your idea, the model's, or forced by an event? Which guardrail do you consider
the single most important, and which one did you add too late?}

\section{How the prompting developed}\label{sec:prompting}

\todo{This section is written from the interview and from the user turns preserved in the session transcripts;
the repository does not contain prompts. Reconstruct the trajectory in four stages and quote the actual
prompts where the author agrees.}

\subsection{Stage 1: build and test (8 September)}
\interview{5}{Your first prompts: how did you describe the tool and the thesis? Did you ask for the tests, or did
the model propose them? What did you do when the tests refuted the thesis?}

\subsection{Stage 2: the research register (9--10 September)}
The commit messages of these two days show a fixed vocabulary appearing: ``reviewer round'', ``rigorous pass'',
``self-review'', ``literature round $n$ (verified)'', ``STATUS'', ``ERRATA''. \evidence{\texttt{git log} 9--10 Sep}
\interview{6}{What did you ask for that produced this? Were there prompts you stopped using because they produced
confident nonsense? How did you learn to ask ``be exact'' (12--13 Sep) and what did that change?}

\subsection{Stage 3: correction under external pressure (12--13 September)}
The author's turns preserved from this phase include the decision principle (``we rather leave it open, or, if we
have reason to think we can fix it, fix it''), the instruction to consider ``the series in totality'', and the
economic constraint (``take care for tokens, maybe use a worse model for a simpler task'').
\interview{7}{How did you decide between the repair routes for paper~IV without being able to check the
mathematics yourself? What did you rely on: the data, the model's account, the structure of the argument?}

\subsection{Stage 4: knowing when to stop (13--14 September)}
The author asked when the internal phase should be considered finished; the answer became the stopping
criterion in the knowledge base. \evidence{\file{KNOWLEDGE.md} \S0b}
\interview{8}{Why did you ask? What did you observe in the model's behaviour that made you suspect over-polishing?}

\section{The error record}\label{sec:errors}

The knowledge base records \NFallacies{} fallacies and the errata \NErrata{} items. We classify them by type and by who
caught them. \todo{Fill the table from \file{KNOWLEDGE.md} \S3 and \file{ERRATA.md}; the classification below is
provisional.}

\begin{center}\small
\begin{tabular}{@{}lp{7.2cm}l@{}}
\toprule
class & examples & caught by\\
\midrule
Borrowed theorem outside its hypotheses & DFI range (F11, F13); Kurokawa's class (KB \S5); Selberg--Delange
non-vanishing (F37) & literature check; outside reader\\
Over-claiming in prose & ``proved'' for a conjecture (F7); $o(1)$ then $O(1)$ for a sum with no bounded remainder
(F15, F26); ``submitted'' when it had not been (F21) & external assessment; self\\
Algebraic slips with structural consequences & pole orders in a ``same proof'' generalisation (F27); orbit
invariance of a coprimality condition (F28); Sali\'e substitution (ERRATA 25) & external assessment; fresh readers\\
Numerics that cannot see what they test & log-term of size $0.002$ against noise $20$ (ERRATA 21); a check that
could not detect $\log\log$ growth & self, after the fact\\
Misjudged effort & ``the remainder bound is a day of work'' (F33) & the work itself\\
Convergence and mode-of-convergence errors & absolute vs.\ paired convergence; integer-$Y$ lemma used at real $Y$
(F35); midpoint of a Fourier series at integers (ERRATA 25) & fresh readers\\
\bottomrule
\end{tabular}
\end{center}

\paragraph{Who caught what.} Internal readings (the model re-reading its own work in a fresh instance) caught
algebraic slips, convergence errors and inconsistencies between sections, reliably and in quantity. The two outside
readers caught what internal readings did not: a claim that is true in the special case and false in general (the
assessment's finding on paper~II), a sum rewritten as an orbit sum over a set that is not an orbit (paper~IV), and a
method invoked with a hypothesis it does not have (paper~I, ERRATA~26). The second outside reader's own assessment of
this asymmetry is quoted in \S\ref{sec:limits}. \evidence{ERRATA 11--15, 18, 22, 24--26}

\paragraph{Two retractions.} The programme retracted two statements it had believed: the sharp form of
Conjecture~1 (there is no bounded remainder at all, F26) and the level-one form of paper~IV's theorem (the object
is a signed sum over $\Gamma_0(e)$-orbits, not one orbit, F28). Both are recorded with the reasoning that had
produced them.
\interview{9}{Your worst moment. Which error did you feel most, and what did you change afterwards?}

\section{Verification}\label{sec:verification}

What ``verified'' means in this programme, from the checklist in the knowledge base (\S6):
\begin{itemize}[itemsep=1pt]
\item Every exact constant by two independent routes.
\item Every explicit formula by exact subtraction, at most one fitted coefficient, and comparison with shifted
  controls (a fit that also matches a wrong frequency proves nothing).
\item Every theorem by at least one adversarial reading by an instance that did not write it.
\item Every prediction with no free parameter tested against data before it is called confirmed: paper~IV's
  smooth-window theorem predicted the amplitude and phase of three spectral lines on two objects from LMFDB
  coefficients alone, and the data matched to a few percent and a few hundredths of a radian (ERRATA 21--23).
\end{itemize}
\interview{10}{What did you personally verify, and how? What could you not verify, and how did you decide to
trust it anyway? Did you ever refuse a claim the model insisted on?}

\section{What the process cannot do}\label{sec:limits}

The second outside reader, after finding the paper~I defect, wrote an assessment of the limits of such readings
that we adopt as the honest statement of the limits of this whole process. Paraphrased: the attacks a model
reader can mount are algebra it can re-derive, numerics it can re-run, and consistency it can grep for; the
failure mode that matters most, a borrowed theorem invoked outside its hypotheses, is detectable only by someone
who has read the sources; its own finding was possible only because Dedekind's conjecture and Siegel zeros are
famous; Shiu's exact uniformity conditions and Katok--Sarnak's normalisation are not, and neither the reader nor the
author can check them. \todo{Quote verbatim with the author's permission.}

Consequently, the state of the mathematics is: every statement has survived at least two adversarial passes and,
for the core of papers~I--II, one external assessment; the class of error that remains possible is known and
named; the only remaining move is review by a specialist who knows the sources. This paper does not claim more.

\section{Recommendations}\label{sec:recommend}

\todo{Write after the interview. Provisional list.}
\begin{enumerate}[itemsep=1pt]
\item Write the thesis as a test suite before writing anything else; let it fail.
\item Keep one file that says what is proved, what is open, and what may not be claimed; make the model read it
  first in every session.
\item Record every error with the rule that would have prevented it; the list is worth more than the theorems.
\item Separate generator and checker: fresh instances, no session context, the papers only.
\item Archive external reviews against the revision hash; apply findings additively in a dated errata file.
\item Decide the repair principle before the repair: correct, or leave open; never a weaker claim about a
  different object.
\item Set the stopping criterion in advance, in writing.
\item Know the class of error the process cannot catch, say so in the papers, and route the work to someone who
  can.
\end{enumerate}

\appendix
\section{The interview}\label{app:interview}
\todo{Verbatim answers to Questions 1--12 (\file{INTERVIEW.md}), lightly edited for length, approved by the
author.}

\section{The process record}\label{app:record}
Generated by \file{scripts/process-record.py} from the repository at the revision named in the footer of this
document. \todo{Add the revision hash and the list of adversarial readings with dates, models and outcomes (from
\file{ERRATA.md} items 18, 22, 24, 25 and \file{KNOWLEDGE.md} \S0c).}

\section*{AI-assisted research: disclosure}
This case study was drafted by the same model that did the work it describes, from the repository record and an
interview with the author; the author verified every factual statement against the record and is responsible for
the text. \todo{Align with the disclosure paragraph of the four papers.}

\end{document}
